% Nouveau EGA V (cote 25): archive cote 25, site page 169 (PDF page 170).
% English translation of the checked French annotated master.
% AI-translated by GPT-6 Astra (OpenAI, ChatGPT Pro mode).
\scanpage{25/169}
\mspage{\hand{\unsure{4}}}

inverse image of the zero section of the vector bundle $V(\hand{\underline{\Omega}}^m_{X/Y}\hand{\otimes}\underline{O}_{V^o})^{\hand{\vee}}\hand{\otimes}\hand{\bigwedge}^m\underline{N}\hand{\otimes}\underline{O}_{V^o})$
under the section corresponding to the homomorphism $(*)_{m}$ ). Since the first member
of $(*_m)$ is of finite type, it also follows that $V^1$ is of finite \unsure{presentation}
over $\unsure{V|u}^o$ . Similarly, $V''$ is defined as the prescheme of zeros of
\[ (*_{m-1}) \qquad \hand{\bigwedge}^{m-1}\underline{N}\hand{\otimes}\underline{O}_{V^1} \hand{\longrightarrow} \hand{\underline{\Omega}}^{m-1}_{X/Y}\hand{\otimes}\underline{O}_{V^1} \]
obtained from $(*)$ by passing to the m-1 th exterior power and restricting to
$V^1$ . As for the invariance \unsure{of the} subprescheme $V^2\hand{\subset} V^1-V''$ , which I have not really
verified and can therefore only conjecture, I think one should be able
to make this point precise as follows. The homomorphism
\[ (*\,*^1) \qquad \underline{E}\hand{\otimes}\underline{O}_{V^1}\hand{\longrightarrow}\hand{\underline{\Omega}}^1_{X/Y}\hand{\otimes}\underline{O}_{V^1} \]
obtained from $(*^o)$ by \var{\xout{restricting to $V^1$\unsure{,} and}}{}{typewriter strike-out}{ega25_25_169-v0001} by composing with the
\uline{surjective} homomorphism $\underline{E}\hand{\otimes}\underline{O}_{V^o}\hand{\to}\underline{N}$ and \var{\xout{composing}}{}{typewriter strike-out}{ega25_25_169-v0002} restricting to $V^1$ , has the same cokernel
$\underline{W}$ as
\[ (*^1) \qquad \underline{N}\hand{\otimes}\underline{O}_{V^1}\hand{\to}\hand{\underline{\Omega}}^1_{X/Y}\hand{\otimes}\underline{O}_{V^1} \quad , \]
which \var{\inserted{$\underline{W}$}}{$\underline{W}$}{insertion}{ega25_25_169-v0003} is therefore locally free over $V^1-V''$ , \hand{(}of rank \var{\over{d}{d}}{d}{writing over}{ega25_25_169-v0004}\var{\deleted{i}}{}{deletion}{ega25_25_169-v0005}\unsure{\var{\deleted{m}}{}{deletion}{ega25_25_169-v0006}|m}+1 at a point
x of $V^1-V''$ at which X has rel dim \hand{d} over Y\hand{)}. It follows that \hand{``}the formation
of the kernel of $(*^1)$ \var{\over{c}{C}}{C}{writing over}{ega25_25_169-v0007}ommutes with every base change $T\hand{\to} V^1$\var{\inserted{$-V''$}}{$-V''$}{insertion}{ega25_25_169-v0008}\hand{''}\var{\over{.}{,}}{,}{writing over}{ega25_25_169-v0009} \var{\over{l}{L}}{L}{writing over}{ega25_25_169-v0010}et $\underline{k}$ be the
kernel of $(*^1)$. If $\underline{K}$ is the kernel of $(*\,*^1)$ , which is therefore invertible over
$V^1-V''$ , one sees \var{\xout{therefore}}{}{typewriter strike-out}{ega25_25_169-v0011} that $\underline{k}$ is a quotient of $\underline{K}$ . On the other hand, on
page 15 we defined a homomorphism
\[ \hand{\underline{M}} : \ \underline{K}\hand{\longrightarrow}\underline{\mathrm{Sym}}^2(\underline{W}) \]
over $V^1-V''$ (it may, be convenient to use $\underline{K}$,$\underline{W}$ exclusively for the Modules
over $V^1-V''$ ...), \var{\deleted{and}}{}{deletion}{ega25_25_169-v0012} \var{\over{t}{T}}{T}{writing over}{ega25_25_169-v0013}he point one would then like to verify
is that this homomorphism in fact factors through the quotient $\underline{k}$ of $\underline{K}$:
\[ \hand{\underline{m}} : \ \underline{k}\hand{\longrightarrow}\underline{\mathrm{Sym}}^2(\underline{W}) \ \var{\xout{\illegible{1}}}{}{typewriter strike-out}{ega25_25_169-v0014} ; \]
and that this latter homomor\unsure{p|b}hism (between two Modules whose definition
\var{\xout{if this point is verified, then}}{}{typewriter strike-out}{ega25_25_169-v0015} no longer depends on u) is in fact independent of
u. If this point is verified (and \uline{if it is true}, an elementary calculation must

\mspage{\hand{169}}
\tnote{25-169-1}{The doubt in presentation is the source's a/e in présentation; of the retains its a/u doubt in du. The repeated by preserves the two live occurrences of en around the canceled clause. The doubtful rank expression is retained without supplying a missing sign. The initial-letter changes in Soit and Le are rendered by the corresponding case changes in Let and The. The extra parenthesis in the first formula and the comma in may, be remain.}
